\documentclass{article}
\usepackage[UTF8]{ctex}
\usepackage{amsmath,amssymb,amsfonts}
\usepackage{enumitem}
\usepackage{geometry}
\geometry{a4paper, margin=1in}
% === 错题本排版优化（不参与题目标识） ===
\setlength{\parindent}{0pt}
\setlength{\parskip}{0.35em plus 0.10em minus 0.08em}
\setlength{\emergencystretch}{3em}
\setlist{itemsep=0.25em, topsep=0.40em, parsep=0pt}
\allowdisplaybreaks[2]
% === 错题本排版优化结束 ===
\begin{document}
\section*{山东大学2025年数学分析真题}
\section*{122．山东大学2025年数学分析试卷}

\subsection{题目1}
1．设 \(g(x)\) 在 \([0,1]\) 上二阶连续可微，满足

\[
g(0)=g^{\prime}(0)=0, g^{\prime \prime}(0)=2 a .
\]

设 \(f(x)=\frac{g(x)}{x}(x \neq 0), f(0)=0\) ，求 \(f^{\prime}(0)\) ．

\subsection{题目2}
2 ．求积分 \(\int_{0}^{1} \frac{\ln (1+x)}{1+x^{2}} \mathrm{~d} x\) ．

\subsection{题目3}
3．若 \(F(t)=\iiint_{x^{2}+y^{2}+z^{2} \leq t^{2}} f\left(x^{2}+y^{2}+z^{2}\right) \mathrm{d} x \mathrm{~d} y \mathrm{~d} z, t>0\) ．其中 \(f\) 是 \(\mathbb{R}\) 上的可微函数．

\begin{itemize}
  \item[（i）] 计算 \(F^{\prime}(t)\) ．
  \item[（ii）] 若 \(f^{\prime}(0)=0\) ，计算 \(\lim _{t \rightarrow 0^{+}} \frac{F(t)}{t^{5}}\) ．
\end{itemize}

\subsection{题目4}
1．设 \(f(x)\) 在 \((0,1)\) 上可导，\(f^{\prime}(x)\) 在 \((0,1)\) 上有界，证明：\(f(x)\) 在 \((0,1)\) 上有界且一致连续．

\subsection{题目5}
2．设 \(f(x, y)=\varphi(|x y|)\) ，且在 \(u=0\) 的附近满足 \(|\varphi(u)| \leq u^{2}\) 。证明：\(f(x, y)\) 在 \((0,0)\) 处可微。

\subsection{题目6}
3．设非负连续函数 \(f(x)\) 在 \([0,+\infty)\) 上单调递减，证明：

\[
F(x)=\int_{0}^{x}(x-2 t) f(t) \mathrm{d} t
\]

是 \([0,+\infty)\) 上的凸函数．

\subsection{题目7}
1．设 \(f(x)=\frac{1}{1-x^{2}}\) ．

\begin{itemize}
  \item[（1）] 求极限 \(\lim _{n \rightarrow \infty}\left[f\left(1-\frac{1}{n+1}\right)-f\left(1-\frac{1}{n}\right)\right]\) 。
  \item[（2）] 证明：\(f(x)\) 在 \([0,1)\) 上不一致连续．
\end{itemize}

\subsection{题目8}
2．设

\[
u_{n}(x)=x^{n} \ln x, x \in(0,1] .
\]

讨论 \(\sum_{n=1}^{\infty} u_{n}(x)\) 在（0 ，1］上的收敛性与一致收敛性，并计算 \(\int_{0}^{1} \sum_{n=1}^{\infty} u_{n}(x) \mathrm{d} x\) ．

\subsection{题目9}
3．设 \(u(x, y, z)\) 是 \(\mathbb{R}^{3}\) 上的连续函数，而且在点 \(\left(x_{0}, y_{0}, z_{0}\right)\) 的某个邻域上有连续的二阶偏导数。记 \(\sum\)是以点 \(\left(x_{0}, y_{0}, z_{0}\right)\) 为球心，\(R\) 为半径的球面，令 \(T(R)=\frac{1}{4 \pi R^{2}} \iint_{\sum} u(x, y, z) \mathrm{d} S, R>0\) ．证明：

\begin{itemize}
  \item[（1）] \(\lim _{R \rightarrow 0^{+}} T(R)=u\left(x_{0}, y_{0}, z_{0}\right)\) ．
  \item[（2）] \(\lim _{R \rightarrow 0^{+}} \frac{T(R)-u\left(x_{0}, y_{0}, z_{0}\right)}{R^{2}}=\frac{1}{6} \Delta u\left(x_{0}, y_{0}, z_{0}\right)\) ，其中 \(\Delta u=\frac{\partial^{2} u}{\partial x^{2}}+\frac{\partial^{2} u}{\partial y^{2}}+\frac{\partial^{2} u}{\partial z^{2}}\) ．
\end{itemize}
\end{document}
